% cuadro C13_3
\begin{table}[htbp]
\centering
\scriptsize
\caption{Techos de endeudamiento, y lo que no es un techo}
\label{cua:C13_3}
\begin{tabular}{p{7.0cm}p{2.6cm}rl}
\toprule
\textbf{Concepto} & \textbf{Fuente} & \textbf{Monto} & \textbf{Unidad} \\
\midrule
Gobierno Federal, endeudamiento neto interno & ILIF art. 2o. & 1 780 000,0 & mdp \\
Gobierno Federal, endeudamiento neto externo & ILIF art. 2o. & 15 500,0 & mdd \\
Petroleos Mexicanos, interno & ILIF art. 4o. & 160 619,6 & mdp \\
Petroleos Mexicanos, externo & ILIF art. 4o. & 5 342,1 & mdd \\
Comision Federal de Electricidad, interno & ILIF art. 4o. & 8 764,2 & mdp \\
Comision Federal de Electricidad, externo & ILIF art. 4o. & 969,0 & mdd \\
Ciudad de Mexico & ILIF art. 6o. & 3 500,0 & mdp \\
--- para contraste, NO son techos --- & --- & --- & --- \\
Deficit presupuestario (DEC art. 2 / CGPE II.6) & CGPE 2026 & 1 393 770,6 & mdp \\
Ingresos derivados de financiamientos (informativo) & ILIF art. 1o. numeral 0 & 1 472 626,4 & mdp \\
\bottomrule
\end{tabular}
\\[2pt]\begin{minipage}{0.92\textwidth}\scriptsize Techo de endeudamiento, déficit y endeudamiento informativo son tres objetos distintos. El techo interno del Gobierno Federal supera al déficit presupuestario en 27.7 por ciento. Fuente: ILIF 2026, artículos 2o., 4o. y 6o.; CGPE 2026. Elaborado por el ITED.\end{minipage}
\end{table}
